\section{Counting, sampling, and the limits of local positivity}
\label{sec:frontier-counting}

Families 113--115 form an integrated case study in converting structural
inequalities into algorithms. Their common question is how to count or sample
a constrained family when positive weights and feasible local moves alone give
little quantitative information. Three different missing implications are
addressed: an energy bound for a restricted class of functions must control
the functions used by an algorithm; an approximate law must have a tractable
exact description before it can be corrected; and a polynomial operation count
must remain polynomial when integer data are written in binary.
A fully polynomial randomized approximation scheme (FPRAS) estimates a
count to relative error $\varepsilon$ with high probability, in time
polynomial in input size and $1/\varepsilon$. The question here is which
new mathematical estimate would make such guarantees possible for the
unrestricted families.

\subsection{General perfect matchings: enlarge the space, then amplify the inequality}

The older matching literature supplies both a method and a precise obstacle.
Jerrum and Sinclair established approximation results through matching walks
\cite{FrontJS1989}; Jerrum, Sinclair and Vigoda obtained the permanent FPRAS by
balancing hole-pattern weights \cite{FrontJSV2004}. The nonbipartite extension
requires additional mathematics: \v{S}tefankovi\v{c}, Vigoda and Wilmes show
torpid mixing for the JSV-type chain under arbitrary hole-pattern weighting
\cite{FrontSVW2017}. Their obstruction concerns that specified state space
and transition mechanism. It is therefore compatible with an algorithm that
changes both.

The release's Theorem 1.1 covers finite simple undirected graphs, rational
\(0<\varepsilon<1\), and \(0<\delta<1/2\), with relative-error failure at most
\(\delta\), exact zero output on empty instances, and polynomial bit time on
every random tape \cite{FrontOpenAIMatching}. Its key sequence is Theorem 3.1's
hole-capacity bound, a weighted subdivision and colored tree-bag enlargement,
and a sampler with local switches and exchanges between two matching
coordinates. Theorem 6.1 bounds additive pair functions, rather than arbitrary
pair functions. Lemma 7.2 supplies the missing upgrade: ordered product
decomposition assigns each interaction residual to at most one coordinate
pair, while replicated adjacent weight tiers make their accumulated energy
absorbable. The replication count is \(32L\), where \(L\) is the large
polynomial comparison bound carried from the pair estimate.
The manuscript derives its claimed polynomial-complexity bound using
deliberately large parameters.

The distinctive review target is this change of analytic scale. A signed
cycle identity is a comparison device, not a path of legal Markov moves;
the full product-space inequality requires its own argument. Reading a
two-coordinate estimate as full mixing would discard the step that makes
the proposed construction unusual. The scope document lists the graph FPRAS,
but only selected accompanying entropy and face statements
\cite{FrontScope113}. Thus the algorithmic theorem and every geometric companion
should retain their separate evidence status.

\subsection{Common bases: restricted observables can be enough}

A matroid abstracts the exchange properties of linear independence.
Its bases are maximal independent sets; a common base must be a base of
each of two matroids on the same ground set. An independence oracle
answers whether a specified set is independent. For a single matroid,
high-dimensional walks already yield an independence-oracle
FPRAS \cite{FrontALOV2024}. Earlier common-base approximations included a
deterministic \(2^{O(r)}\) factor \cite{FrontAOV2021}. Their algebraic framework
does not automatically survive intersection. On the four edges of \(K_{2,2}\),
the common bases of its two vertex-partition matroids have generating
polynomial
\[
g(x)=x_1x_4+x_2x_3.
\]
Its Hessian has eigenvalues \(1,1,-1,-1\); consequently it fails the
one-positive-direction condition for a nonnegative quadratic Lorentzian
polynomial. This elementary example isolates a structural obstruction,
rather than a hardness conclusion.

The September manuscript's Theorem 1.1 uses arbitrary equal-rank matroids
on an explicitly enumerated ground set, with exact independence oracles.
It bounds oracle calls and additional bit work polynomially on every
execution \cite{FrontOpenAICommon}. A paired reduction uses
\(M_1\oplus M_2^*\), then positive rank-deficiency weights and one-empty/one-full
pair defects. Br\"and\'en--Huh quadratic signatures \cite{FrontBH2020} control
a recursively balanced signed current. Lemma 4.1 bounds transport between
missing-slot conditional laws. Proposition 5.3 controls transversal values
and defect-class means, leaving within-defect fluctuations uncontrolled.
Lemma 6.2 separately bounds the transversal trace. Fresh warm restarts and
observable variance estimates then support annealing and learned defect
multipliers. The scope document specifies the independence-oracle matroid
theorem \cite{FrontScope114}; it does not identify the later binary-capacity
integer-polymatroid extension as part of that coverage.

Chen, Vigoda, and Yang already use partition-restricted Poincar\'e
control to estimate hole-pattern frequencies for the permanent
\cite[Definition~4 and Theorem~5(ii)]{FrontCVY2026}. Their framework
bounds stationary empirical averages of block-constant observables
under the specified reversible-chain hypotheses. The common-base
manuscript explicitly adopts this framework. Its candidate advance is
the matroid-specific transport estimate together with compatible
transversal trace and warm-start control. These additional estimates
must connect the restricted inequality to the observables and returns
used by the counting algorithm.

A contemporaneous comparison is also essential:
Chen and Liu's 5 October 2026 preprint also gives a common-base FPRAS
and develops a nonnegative-curvature framework for Hadamard products of
measures \cite{FrontChenLiu2026}. Its field-selection and curvature mechanism
provides a separate comparison route. Manuscript dates alone therefore
do not establish priority, and the September construction should be assessed
through its specific transport and simulation arguments.

\subsection{Contingency tables: make exactness affordable in expectation}

The target law here is uniform over nonnegative integer tables with prescribed
margins. Dense-margin sampling has longstanding geometric predecessors
\cite{FrontDKM1997,FrontMorris2002}, while perfect sampling was already available
for two-row tables \cite{FrontKijimaMatsui2003}. These restrictions matter:
a result for fixed row count or sufficiently large margins does not give
a polynomial bound jointly in both dimensions and the margin bit length.

The release's Theorem 1.1 allows arbitrary nonnegative equal-total margins,
both dimensions variable, and no cell restrictions. It supplies error
\(2^{-k}\) in worst-case polynomial bit time and exact uniform sampling in
expected polynomial bit time \cite{FrontOpenAITables}. Small-margin cells
form explicit defect states; padding makes the remaining rectangular
completion problems dense. Conditional transport and repair lead to the
small-state Poincar\'e inequality in Theorem 4.3. The exactness bottleneck
is Proposition 6.3: the actual bounded sampler's law, including fallbacks,
is tabulated with an exponential input-size factor but polynomial dependence
on precision. Theorem 6.4 applies the earlier residual-mixture construction
of G\"obel, Liu, Manurangsi and Pappik \cite{FrontGLMP2024}.
Their method requires usable output-law computation, not merely a small
total-variation error.

For any proposed exact correction, write \(\delta\) for the actual
correction-branch probability and \(T_{\rm ordinary}\) for the mean work
conditioned on the ordinary branch. The decisive accounting is
\[
\mathbb E[T]=(1-\delta)T_{\rm ordinary}
+\delta\,\mathbb E[T_{\rm correction}].
\]
The correction branch may be expensive while its contribution to the
expectation remains controlled. An implementation that instead enumerates
all precision-dependent random tapes can destroy that accounting. This
distinguishes a probability identity from an efficient sampler.
The scope document likewise separates unrestricted exact sampling from
cell-bounded approximate counting \cite{FrontScope115}; forbidden cells
cannot simply be inserted into the unrestricted sampler's guarantee.

\subsection{From algorithmic guarantees to implementations}

The three constructions repair different interfaces between analysis and
computation. Matching replication upgrades a restricted inequality;
common-base trace simulation exploits a restricted inequality without
that upgrade; table-law tabulation makes a rare exact correction computable.
Calling all three \emph{rapid mixing from positivity} would obscure the
specific work.

A useful research program would isolate these interfaces as reusable
lemmas with independently checked hypotheses. For common bases, one could
compare conditional transport and curvature on the same oracle families,
measuring which quantities each certificate controls. For tables, the
natural implementation task is a finite-kernel law evaluator whose state
enumeration is separated from its requested precision. For matchings, a
small-instance prototype should report enlarged graph size, replication,
and acceptance probabilities before extrapolating runtime. These are
proposals motivated by the proof structures, not performance conclusions.

The theorem statements concern explicitly enumerated graphs or ground sets,
exact oracles, and binary integer input. They provide an algorithmic frontier
if the proposed proofs hold. Practical utility additionally requires feasible
constants, a faithful target distribution, and measured cost against existing
restricted-family algorithms. The scope documents identify the supplied
formalization targets; assessing the supplied proof architecture and
running a proof checker are distinct verification activities.

\section{Unique Games and the completeness barrier}
\label{sec:frontier-unique-games}

A Unique Game assigns labels to vertices, with each edge requiring
one label to be a specified bijective function of the other. The goal
is to satisfy as many edges as possible. Completeness describes the
fraction achievable in the yes instances of a reduction; soundness
bounds the fraction achievable in its no instances. Family 102 claims
hardness even when these fractions are near one and near zero,
respectively. The established shortcode and Grassmann-expansion line gives
imperfect-completeness hardness \cite{FrontBKS2019,FrontKMS2023}. Near-one
completeness is a separate requirement. Ordinary repetition multiplies
completeness as well as reducing soundness, so it does not itself repair a
one-half completeness barrier. Khot's survey states the independent-error
quantifiers of Unique Games \cite{FrontKhot2010}, and Raghavendra's
SDP-optimality framework gives consequences conditional on that hypothesis
\cite{FrontRaghavendra2008}.

The release's Theorem 1.1 asserts a deterministic polynomial-time reduction
from 3SAT, for every fixed \(0<\varepsilon,\delta<1/2\), to simple unweighted
bipartite translation games over \(\mathbb F_2^s\), where
\(s=s(\varepsilon,\delta)\) depends only on the fixed error parameters,
with completeness
at least \(1-\varepsilon\) and soundness at most \(\delta\)
\cite{FrontOpenAIUGC}. Lemma 3.1 proposes a latent-alphabet gadget: an equivariant
nonlinear observable is stable under noise, while families of linear
observations with sufficiently high restriction rank on the gadget's
distinguished output subspace \(K\) detect that noise with probability at least
\(1/8\). Section 5 uses Fourier comparison to return to the ordinary
degree-two shortcode inverse theorem. A correlated-advice decoder and
clean-coordinate repetition supply the soundness bridge.
The scope document describes the translation-game theorem and separately
lists Max-Cut, cover and deletion reductions \cite{FrontScope102}.

The gadget seeks to keep a nonlinear signal stable while making a
sufficiently rich family of linear probes detect the perturbation.
This is the proposed way to preserve high completeness without losing
the soundness argument. The new object is that simultaneous separation,
while the consequences for SDP thresholds have an established conditional
theory. Its evaluation should concentrate on uniformity over adaptive
linear lifts, the row-fiber transfer, and the information actually visible to
each decoded strategy. H{\aa}stad's earlier Fourier decoding explains why
extracting a strategy need not mean recovering a nearby codeword
\cite{FrontHastad2001}; the new gadget must still preserve the particular
completeness and soundness required here.

The distinction from a routine amplification step is thus mathematical:
ordinary repetition reduces both completeness and soundness, whereas
this gadget is intended to separate their behavior. That gain depends
on stability under adaptive linear lifts and on a decoder that uses
only the information available to a valid strategy.
